% !TEX TS-program = pdflatexmk
\RequirePackage{iftex}
\ifPDFTeX
  \pdfoutput=1
  \pdfminorversion=7
  \input glyphtounicode
  \pdfgentounicode=1
\fi
\ifdefined\pdfinterwordspaceon\pdfinterwordspaceon\fi
\documentclass[justified,notoc,twoside,nobib,nofonts]{tufte-handout}
\usepackage[T1]{fontenc}
\usepackage[utf8]{inputenc}
\usepackage[english]{babel}
\ifPDFTeX
  \usepackage[tracking]{microtype}
\else
  \usepackage{microtype}
\fi
\usepackage{amsthm,thmtools,thm-restate}
\usepackage{amsmath,amssymb,mathtools}
% KPfonts text, with the same New PX maths as the font comparison.
\usepackage[nomath,oldstylenums]{kpfonts}
\edef\KPbodyfamily{\rmdefault}
\renewcommand{\rmdefault}{zpltlf}
\usepackage{newpxmath}
\let\rmdefault\KPbodyfamily
\renewcommand{\familydefault}{\rmdefault}

% BaskervilleF headings use native Type 1 fonts and work with pdfLaTeX.
\newcommand{\BaskervilleHeading}{%
  \fontencoding{T1}\fontfamily{BaskervilleF-LF}%
  \fontseries{m}\fontshape{n}\selectfont}
\titleformat{\section}[hang]
  {\BaskervilleHeading\fontsize{15}{18}\selectfont}
  {\thesection}{1em}{}
\titleformat{\subsection}[hang]
  {\BaskervilleHeading\fontsize{13}{16}\selectfont}
  {\thesubsection}{1em}{}
\titleformat{\paragraph}[runin]
  {\BaskervilleHeading}{\theparagraph}{1em}{}
\usepackage{etoolbox}
\makeatletter
\patchcmd{\maketitle}{\LARGE\itshape\@title}
  {\BaskervilleHeading\fontsize{22}{26}\selectfont\raggedright\@title}
  {}{\PackageError{defect-three-oriented-cdc}{Title font patch failed}{}}
% The author and email are printed immediately below the title block.
\patchcmd{\maketitle}{\@@author\@@date}{\@@date}
  {}{\PackageError{defect-three-oriented-cdc}{Author block patch failed}{}}
\makeatother
\usepackage{booktabs,array,enumitem,needspace}
\usepackage[numbers,square]{natbib}
\usepackage{xcolor}
\usepackage{xurl,hyperref,cleveref}
\urlstyle{same}
\geometry{a4paper,top=23mm,bottom=46mm,left=28mm,right=56mm,marginparwidth=3cm}
\setlength{\emergencystretch}{2em}
\clubpenalty=10000
\widowpenalty=10000
\displaywidowpenalty=10000
\setcounter{secnumdepth}{2}
\AtBeginDocument{\fontsize{12}{15}\selectfont}
\fancyhead[LE]{\thepage\quad\textsmallcaps{nikolay ulyanov}}
\fancyhead[RO]{\textsmallcaps{oriented five-cycle covers}\quad\thepage}
\declaretheorem[name=Theorem,numberwithin=section]{theorem}
\declaretheorem[name=Proposition,sibling=theorem]{proposition}
\declaretheorem[name=Lemma,sibling=theorem]{lemma}
\declaretheorem[name=Corollary,sibling=theorem]{corollary}
\declaretheorem[name=Definition,sibling=theorem,style=definition]{definition}
\setlist[enumerate]{label=\textup{(\roman*)},leftmargin=*,itemsep=2pt,topsep=4pt}
\crefname{theorem}{theorem}{theorems}
\crefname{lemma}{lemma}{lemmas}
\crefname{definition}{definition}{definitions}
\setcitestyle{numbers,square}

% Full references appear in the margin at their first citation.
% Margin selection uses the canonical latex-margin-citations skill:
% margin notes stay visible but do not enter the main copied text.
\makeatletter
\newdimen\@gp@fnmarkwidth
\newcommand{\bothcite}[2][0cm]{%
  \citep{#2}\sidenote[][#1]{\raggedright\csname paper@#2\endcsname}%
  \kern\@gp@fnmarkwidth\@gp@cite@space}
\def\@gp@cite@space{\@ifnextchar.{\@gp@cite@halfspace}{\@gp@cite@comma}}
\def\@gp@cite@comma{\@ifnextchar,{\@gp@cite@halfspace}{\@gp@cite@semicolon}}
\def\@gp@cite@semicolon{\@ifnextchar;{\@gp@cite@halfspace}{\@gp@cite@colon}}
\def\@gp@cite@colon{\@ifnextchar:{\@gp@cite@halfspace}{\@gp@cite@exclamation}}
\def\@gp@cite@exclamation{\@ifnextchar!{\@gp@cite@halfspace}{\@gp@cite@question}}
\def\@gp@cite@question{\@ifnextchar?{\@gp@cite@halfspace}{\@gp@cite@openparen}}
\def\@gp@cite@openparen{\@ifnextchar({\@gp@cite@halfspace}{\@gp@cite@closeparen}}
\def\@gp@cite@closeparen{\@ifnextchar){\@gp@cite@halfspace}{\@gp@cite@bracket}}
\def\@gp@cite@bracket{\@ifnextchar]{\@gp@cite@halfspace}{\@gp@cite@normalspace}}
\def\@gp@cite@halfspace{\nobreak\hskip.5\fontdimen2\font\relax}
\def\@gp@cite@normalspace{\nobreak\space}
% GP margin selection v1 (pdfLaTeX). Keep this block in the standalone source.
\makeatletter
\@ifundefined{@gp@fnmarkwidth}{\newdimen\@gp@fnmarkwidth}{}
\ifPDFTeX
\def\@makefnmark{%
  \hbox{%
    \setbox\z@=\hbox{\@textsuperscript{\normalfont\footnotesize\@thefnmark}}%
    \global\@gp@fnmarkwidth=\wd\z@%
    \setbox\tw@=\copy\z@%
    \immediate\pdfxform\tw@%
    \pdfannot width\wd\z@ height\ht\z@ depth\dp\z@{%
      /Subtype /Watermark /F 708 /Contents ()%
      /AP << /N \the\pdflastxform\space 0 R >>%
    }%
  }%
}
\fi
% Selection-safe margin text: annotation appearances, with separate URI links.
\ifPDFTeX
\newbox\@gp@marginbox
\newdimen\@gp@marginwidth
\newdimen\@gp@marginheight
\newdimen\@gp@marginfirstheight
\newdimen\@gp@marginpadding
\newdimen\@gp@marginskip
\newdimen\@gp@margindepth
\newcount\@gp@marginserial
\long\def\@gp@marginappearance#1{%
  \begingroup
  \global\let\@gp@marginuri\@empty
  \setbox\@gp@marginbox=\vtop{%
    \hsize=\marginparwidth
    \def\href##1##2{%
      \ifx\@gp@marginuri\@empty
        \xdef\@gp@marginuri{\detokenize{##1}}%
      \else
        \PackageError{gp-margin-selection}{More than one URI in a margin note}%
          {Use a separate note for each reference.}%
      \fi
      \textcolor{\@urlcolor}{##2}%
    }%
    #1\par
  }%
  \@gp@marginfirstheight=\ht\@gp@marginbox
  \setbox\@gp@marginbox=\vbox{\unvbox\@gp@marginbox}%
  \@gp@marginpadding=\baselineskip
  \advance\@gp@marginpadding-\prevdepth
  \advance\@gp@marginpadding-\@gp@marginfirstheight
  \ifdim\@gp@marginpadding<\lineskiplimit
    \@gp@marginpadding=\lineskip
  \fi
  \@gp@marginskip=\baselineskip
  \advance\@gp@marginskip-\prevdepth
  \advance\@gp@marginskip-\ht\@gp@marginbox
  \ifdim\@gp@marginskip<\lineskiplimit
    \@gp@marginskip=\lineskip
  \fi
  \advance\@gp@marginpadding-\@gp@marginskip
  \vskip\@gp@marginpadding
  \@gp@marginwidth=\wd\@gp@marginbox
  \@gp@marginheight=\ht\@gp@marginbox
  \@gp@margindepth=\dp\@gp@marginbox
  \immediate\pdfxform\@gp@marginbox
  \edef\@gp@marginform{\the\pdflastxform}%
  \global\advance\@gp@marginserial\@ne
  \hbox{%
    \pdfannot width\@gp@marginwidth height\@gp@marginheight depth\@gp@margindepth{%
      /Subtype /Watermark /F 708 /Contents ()%
      /NM (GPmargin-\the\@gp@marginserial)%
      /AP << /N \@gp@marginform\space 0 R >>%
    }%
    \ifx\@gp@marginuri\@empty\else
      \pdfannot width\@gp@marginwidth height\@gp@marginheight depth\@gp@margindepth{%
        /Subtype /Link /F 4 /Border [0 0 0]%
        /A << /S /URI /URI (\pdfescapestring{\@gp@marginuri}) >>%
      }%
    \fi
    \vrule width0pt height\@gp@marginheight depth\@gp@margindepth
    \kern\@gp@marginwidth
  }%
  \endgroup
}
\else
\long\def\@gp@marginappearance#1{#1}
\fi
\renewcommand\@footnotetext[2][0pt]{%
  \marginpar{%
    \hbox{}\vspace*{#1}%
    \def\baselinestretch{\setspace@singlespace}%
    \reset@font\footnotesize%
    \@tufte@margin@par%
    \vspace*{-1\baselineskip}%
    \@gp@marginappearance{%
      \noindent
      \protected@edef\@currentlabel{%
        \csname p@footnote\endcsname\@thefnmark%
      }%
      \color@begingroup
        \@makefntext{\ignorespaces#2}%
      \color@endgroup
    }%
  }%
}
\renewcommand\marginnote[2][0pt]{%
  \ifthenelse{\boolean{@tufte@loadnatbib}}{%
    \let\cite\@tufte@infootnote@cite%
  }{}%
  \gdef\@tufte@citations{}%
  \marginpar{%
    \hbox{}\vspace*{#1}%
    \@tufte@marginnote@font\@tufte@marginnote@justification%
    \@tufte@margin@par\vspace*{-1\baselineskip}%
    \@gp@marginappearance{\noindent #2}%
  }%
  \@tufte@print@citations%
  \ifthenelse{\boolean{@tufte@loadnatbib}}{%
    \let\cite\@tufte@normal@cite%
  }{}%
}
\makeatother

\expandafter\def\csname paper@KarabasEtAl2024\endcsname{%
  J\'an Karab\'a\v{s}, Edita M\'a\v{c}ajov\'a, Roman Nedela and
  Martin \v{S}koviera. \emph{Cubic graphs with colouring defect 3}.
  Electronic Journal of Combinatorics 31(2) (2024), P2.6.
  \href{https://doi.org/10.37236/12333}{doi:10.37236/12333}.}
\newcommand{\paperreference}[1]{\csname paper@#1\endcsname}
\makeatother

\title[Graph Puzzles III.3-preview: Oriented five-cycle double covers]{Graph Puzzles III.3-preview: Oriented five-cycle double covers of snarks with colouring defect three}
\author{Nikolay Ulyanov + AI-Slop}
\date{30 August 2026}
\hypersetup{
  pdftitle={Graph Puzzles III.3-preview: Oriented five-cycle double covers of snarks with colouring defect three},
  pdfauthor={Nikolay Ulyanov + AI-Slop},
  colorlinks=true,
  linkcolor=blue!50!black,
  citecolor=blue!50!black,
  urlcolor=blue!50!black
}

\begin{document}
\selectlanguage{english}
\maketitle
Nikolay Ulyanov\textsuperscript{*} + AI-Slop%
\marginnote{\textsuperscript{*}{\scriptsize\href{mailto:ulyanick@gmail.com}{ulyanick@gmail.com}}}\par

\begin{abstract}
\noindent\newthought{Abstract: }
For every optimal triple of perfect matchings of a snark of colouring defect
three, the triple's hexagonal core extends to an oriented \(5\)-cycle double
cover.  We construct the five directed cycle members.  Three
auxiliary orientations of the bichromatic
exterior paths determine an ordered pair of cover labels on every edge.
At each cubic vertex the three ordered pairs form a directed triangle.
Consequently every label induces a directed even subgraph, while every edge
belongs to exactly two such subgraphs with opposite directions.
\end{abstract}
\noindent\rule{\linewidth}{0.4pt}

\section{An ordered-pair criterion}

Here a \emph{cycle} means an even subgraph and may be disconnected.
We use \emph{snark} for a connected bridgeless cubic graph with no proper
\(3\)-edge-colouring.  A \emph{\(3\)-array} is a triple of perfect matchings.
The multiplicity of an edge is the number of matchings of the triple that
contain it; the \emph{core} of the triple consists of the edges whose
multiplicity is not one.  The \emph{colouring defect} of a cubic graph is the
minimum, over all \(3\)-arrays, of the number of edges of multiplicity zero.

\begin{definition}\label{def:o5cdc}
An \emph{oriented \(5\)-cycle double cover} of a graph \(G\) is an indexed
family \((D_1,\ldots,D_5)\) of directed cycles such that every edge belongs
to exactly two members and receives opposite directions in those two
members.
\end{definition}

\Needspace{16\baselineskip}
\begin{lemma}[Ordered-pair criterion]\label{lem:pairs}
Give every edge a reference direction and an ordered pair \(p\to q\) of
distinct labels from \(\{1,\ldots,5\}\).  Let \(D_p\) use the edge in its
reference direction and let \(D_q\) use it in the reverse direction.
Suppose that, when the three incident edges at each vertex are read
outwards---retaining \(p\to q\) if the reference direction points out of the
vertex and reversing it otherwise---their ordered pairs form a directed
triangle.  Then
\((D_1,\ldots,D_5)\) is an oriented \(5\)-cycle double cover.
\end{lemma}

\begin{proof}
Each edge lies in precisely the two members named by its pair, with opposite
directions.  At a vertex, every label occurring in the directed triangle
occurs once as a tail and once as a head.  Hence each \(D_r\) has one edge
entering and one leaving the vertex, or no incident edge there.  Thus every
\(D_r\) is a directed even subgraph.
\end{proof}

\section{The defect-three construction}

\begin{theorem}\label{thm:main}
Let \(G\) be a snark with colouring defect \(3\).  For every optimal triple
of perfect matchings of \(G\), there is an oriented \(5\)-cycle double cover
having the triple's induced hexagonal core as one indexed member.
\end{theorem}

\begin{proof}
Fix an arbitrary triple of perfect matchings attaining the colouring defect.
By the
defect-three core characterization
\bothcite{KarabasEtAl2024}, Theorem~3.3, their core is an induced cycle
\[
 C=v_1v_2\cdots v_6v_1.
\]
Every edge of \(G-E(C)\) is singly covered, and membership in the three
matchings gives a proper \(3\)-edge-colouring of \(G-E(C)\).  After rotating
\(C\) and permuting the colours, the six boundary edges at
\(v_1,\ldots,v_6\) have colours
\begin{equation}\label{eq:boundary}
 1,1,2,2,3,3.
\end{equation}
Let \(F_i\) denote the exterior colour-\(i\) class.

\subsection*{Three auxiliary orientations}

Fix \(i\in\{1,2,3\}\), write
\(\{i,j,k\}=\{1,2,3\}\), and put
\[
 P_i=F_j\cup F_k,\qquad K_i=C\cup P_i.
\]
Every vertex outside \(C\) has degree two in \(P_i\).  Exactly four
vertices of \(C\) are incident with \(P_i\), so \(P_i\) consists of two
paths pairing these four terminals, together with disjoint circuits.
Suppressing every degree-two vertex in the component of \(K_i\) containing
\(C\) produces a cubic multigraph \(Q_i\) consisting of a \(4\)-cycle and a
perfect matching.

We claim that \(Q_i\) is bipartite.  By rotating and relabelling, it is
enough to take \(i=1\).  The four terminals are
\(v_3,v_4,v_5,v_6\), in that cyclic order.  A \(4\)-cycle plus a perfect
matching is non-bipartite only when the matching pairs the opposite
vertices \(v_3v_5\) and \(v_4v_6\).  If this occurred, interchange colours
\(2\) and \(3\) on the exterior Kempe path joining \(v_3\) to \(v_5\).
The boundary word would become
\[
 1,1,3,2,2,3.
\]
This colouring extends across \(C\): give the consecutive cycle edges
\(v_rv_{r+1}\), with indices modulo \(6\), the colours
\[
 3,2,1,3,1,2,
\]
respectively.  At every \(v_r\), the boundary edge and the two cycle edges
then receive all three colours.  This would Tait-colour \(G\), a
contradiction.  Hence every \(Q_i\) is bipartite.

Choose a bipartition \(A_i\cup B_i\) of \(Q_i\).  Orient its four
suppressed \(C\)-segments from \(B_i\) to \(A_i\), and orient its two
matching edges from \(A_i\) to \(B_i\).  Lift these directions consistently
through the suppressed paths, and orient every circuit component of
\(P_i\) cyclically.  Call the resulting orientation of \(K_i\)
\(\mathcal O_i\).

Thus \(\mathcal O_i\) passes through every degree-two vertex.  At a core
vertex whose boundary edge has colour \(i\), its two \(C\)-edges point in
opposite directions.  At either endpoint of a \(P_i\)-path, the two
\(C\)-edges both point in and the exterior edge points out, or all three
directions are reversed.

\subsection*{The five cover labels}

Fix an arbitrary reference direction on every edge.  For
\(e\in K_i\), put \(\sigma_i(e)=+\) or \(-\) according as
\(\mathcal O_i\) agrees or disagrees with the reference direction; for
\(e\in F_i\), put \(\sigma_i(e)=0\).  The sign word
\[
 \sigma(e)=(\sigma_1(e),\sigma_2(e),\sigma_3(e))
\]
assigns an ordered pair to \(e\) by the following table:
\begingroup
\small
\begin{equation}\label{eq:pair-table}
\begin{array}{@{}c@{\quad}l@{\quad}l@{}}
\toprule
\text{edge class}&\text{displayed sign words}&\text{ordered pairs}\\
\midrule
C   & +++,\ +--,\ -+-,\ --+&
       1\to5,\ 2\to5,\ 3\to5,\ 4\to5\\
F_1 & 0++,\ 0+-&1\to2,\ 3\to4\\
F_2 & +0+,\ +0-&1\to3,\ 2\to4\\
F_3 & ++0,\ +-0&1\to4,\ 2\to3\\
\bottomrule
\end{array}
\end{equation}
\endgroup
The negative of a displayed sign word receives the reversed ordered pair.
Use the pair \(p\to q\) in \cref{lem:pairs}: the member \(D_p\) takes the
reference direction and \(D_q\) takes the reverse direction.
In the vector calculations below we identify \(+\) and \(-\) with \(+1\)
and \(-1\).  At an endvertex \(v\) of an edge \(e\), let
\(\sigma_v(e)=\sigma(e)\) if the reference direction points away from \(v\),
and let \(\sigma_v(e)=-\sigma(e)\) otherwise.  Thus \(\sigma_v(e)\) is the
sign word read outwards at \(v\), and its table pair is reversed in the
second case.

It remains only to check the local directed-triangle condition.  Put
\[
\begin{aligned}
 t_1&=(1,1,1),&t_2&=(1,-1,-1),\\
 t_3&=(-1,1,-1),&t_4&=(-1,-1,1).
\end{aligned}
\]
The four vectors \(t_1,\ldots,t_4\) are affinely independent.  The first
row of \eqref{eq:pair-table} says that a core sign word is \(t_p\) for
\(p\to5\), or \(-t_p\) for \(5\to p\).  The remaining rows say that an
exterior sign word belonging to the pair \(p\to q\) is
\[
 \frac{t_p-t_q}{2}.
\]

First let \(v\notin V(C)\).  For each coordinate \(i\), the orientation
\(\mathcal O_i\) passes through \(v\) on the two edges whose colours differ
from \(i\).  Their two outward \(i\)-signs are therefore opposite, while the
colour-\(i\) edge has sign \(0\).  Hence the three vectors \(\sigma_v(e)\)
sum to zero.  If their outward-read ordered pairs are
\(p_s\to q_s\), then
\[
 \sum_{s=1}^3(t_{p_s}-t_{q_s})=0.
\]
For \(r\in\{1,2,3,4\}\), let \(d_r\) be the number of these pairs with tail
\(r\), minus the number with head \(r\).  The displayed equality says
\(\sum_r d_rt_r=0\), while automatically \(\sum_r d_r=0\).  Affine
independence gives \(d_r=0\) for every \(r\).  Thus every label occurs
equally often as a tail and as a head.  Three non-loop ordered pairs with
this property form a directed triangle.

Now let \(v\in V(C)\), and let its exterior edge \(a\) have colour \(i\).
Suppose that the outward-read table pair of \(a\) is \(p\to q\).  Then
\[
 \sigma_v(a)=\frac{t_p-t_q}{2},\qquad
 (t_p)_i=(t_q)_i,\qquad
 (t_p)_j=-(t_q)_j\quad(j\ne i).
\]
In \(\mathcal O_i\), the two \(C\)-edges have opposite signs.  In each
\(\mathcal O_j\) with \(j\ne i\), they have the same sign, opposite to that
of \(a\).  Consequently their two sign words are \(t_q\) and \(-t_p\).
The table labels the two \(C\)-edges by \(q\to5\) and \(5\to p\).
Together with \(p\to q\), they form the directed triangle
\[
 p\longrightarrow q\longrightarrow5\longrightarrow p.
\]

The hypothesis of \cref{lem:pairs} now holds at every vertex.  Therefore
\((D_1,\ldots,D_5)\) is an oriented \(5\)-cycle double cover.  By the
table, label \(5\) occurs precisely on the core edges.  Hence
\(E(D_5)=E(C)\), and \(D_5\) is a cyclic orientation of \(C\).

Finally, all five members are nonempty.  Indeed, if some \(D_r\) were
empty, all edge-pairs would lie in the complete graph on the remaining
four labels.  Colour each edge of \(G\) by the unique one of the three
one-factors of this \(K_4\) containing its pair.  At every vertex the
three pairs form a triangle and hence receive the three distinct colours.
This would Tait-colour \(G\), again a contradiction.
\end{proof}

\Needspace{8\baselineskip}
\begingroup
\fontsize{10}{12}\selectfont
\setlength{\bibsep}{5pt}
\begin{thebibliography}{KMNS24}
\bibitem[KMNS24]{KarabasEtAl2024}\paperreference{KarabasEtAl2024}
\end{thebibliography}
\endgroup

\end{document}
